% CS F407 Artificial Intelligence: Final Academic Project Report
% Allowed: 4 double column pages, including references in this TeX template.
\documentclass[11pt,a4paper,twocolumn]{article}

\usepackage[margin=2cm,columnsep=0.8cm]{geometry}
\usepackage[T1]{fontenc}
\usepackage{lmodern}
\usepackage{microtype}
\usepackage{amsmath,amssymb}
\usepackage{graphicx}
\usepackage{booktabs}
\usepackage{tabularx}
\usepackage{enumitem}
\usepackage{titlesec}
\usepackage[hidelinks]{hyperref}

\titleformat{\section}{\normalsize\bfseries}{\thesection}{0.6em}{}
\titlespacing*{\section}{0pt}{1.2ex plus 0.3ex}{0.6ex}

\setlist{nosep,leftmargin=*}
\setlength{\parskip}{0.3ex}

\newcommand{\coursename}{CS F407 Artificial Intelligence, AY 2026--27 Semester I}
\newcommand{\teamno}{Team Thinkers}
\newcommand{\githuburl}{https://github.com/LearnerKg23/Audio-Engineering}
\newcommand{\demourl}{https://learnerkg23.github.io/Audio-Engineering/}

\begin{document}

\twocolumn[{%
\begin{center}
  {\large\bfseries Who Spoke What and When: Modular Speech Disentanglement and Conversational Processing for Overlapping Code-Switched Audio}\\[1ex]
  \teamno\\
  Kushal Gupta (2023B1A70698G), \quad Vansh Sharma (2023B1A70715G),\\[0.3ex]
  Dorothy (2023B1A70714G), \quad Pranav (2024A7PS0515G)\\[0.5ex]
  {\small \coursename, BITS Pilani, K.K. Birla Goa Campus}\\[1.5ex]
  \begin{minipage}{0.85\textwidth}
    \small
    \centerline{\bfseries Abstract}\vspace{0.5ex}
    Overlapping speech impairs speaker diarization and automated transcription, especially in colloquial Hindi-English (Hinglish) code-switched conversations. We evaluate a modular five-stage pipeline: audio conditioning, time-domain separation, spectral-clustering diarization, Indic CTC acoustic modeling, and schema-constrained dialogue structuring with extractive lexical retrieval. In Track~A controlled benchmarks, pretrained Conv-TasNet achieves a mean SI-SDR improvement ($\Delta$SI-SDR) of +15.45~dB (mean output: 15.45~dB, baseline: $-0.00$~dB [$-0.003$~dB exact]) across six two-speaker mixtures. Across nine controlled mixtures, separation reduces mean Diarization Error Rate (DER) from 40.52\% to 23.97\%. On a bilingual mixture, Vakyansh Hindi Wav2Vec2 Word Error Rate (WER) drops from 86.28\% on raw audio to 51.57\% with separation and diarization. On Track~B simulated broadcast and cellular-bandpass stress clips (27.3~s total), evaluation exposes speaker under-counting under dense overlap (58.7\% DER) and narrowband degradation (63.0\% WER).
  \end{minipage}
\end{center}
\vspace{1.5ex}
}]

\section{Introduction}
Overlapping speech in single-channel Indian conversations causes end-to-end automated speech recognition (ASR) models to drop subordinate speakers or collapse into repetitive loops. Intra-sentential code-switching between Hindi and English (``Hinglish'') further mixes Devanagari and Latin vocabulary, disrupting monolingual acoustic and language priors.

The pipeline takes \textbf{input} single-channel overlapping Hindi-English conversational audio (16~kHz PCM WAV) and produces \textbf{output} time-stamped, speaker-attributed transcriptions, Pydantic-structured dialogue summaries, and extractive BM25 conversational retrieval. Rather than training backbones from scratch, the authors contributed unified stage contracts, polyphase resampling (16$\leftrightarrow$8~kHz) around pretrained separation, native 192-d spectral clustering with eigengap estimation, CTC delimiter reconstruction and NFKC normalization, rule-based entity grounding checks, and reproducible Track~A/B evaluation suites. Because benchmarks test distinct conditions, results retain explicit sample descriptors rather than pooled averages.

\section{Background and related work}
Joining source separation, diarization, and Indic CTC modeling is essential because monolithic transformers suffer attention collapse on concurrent Hinglish speech, while standalone Indic CTC models drop overlapping turns and music-domain separators collapse human voices into one stem. Decoupling waveform separation from turn-level acoustic clustering isolates speakers before Devanagari decoding and structured retrieval.

Conv-TasNet~\cite{luo2019convtasnet} uses a 1D convolutional encoder, dilated temporal convolutional network (TCN) separator, and linear decoder via permutation invariant training, with torchaudio weights~\cite{yang2022torchaudio} trained on 8~kHz Libri2Mix~\cite{cosentino2020librimix}. Hybrid Demucs~\cite{defossez2021demucs} pairs time convolutions with spectrogram U-Nets for music demixing rather than speech separation. Speaker diarization extracts windowed embeddings~\cite{bredin2020pyannote} grouped via normalized-Laplacian spectral clustering~\cite{ng2001spectral}. Separation is evaluated by Scale-Invariant Signal-to-Distortion Ratio (SI-SDR)~\cite{leroux2019sisdr} and BSS-Eval SDR~\cite{vincent2006performance}, with $\Delta\text{SI-SDR} = \text{SI-SDR}(\hat{s}, s) - \text{SI-SDR}(y, s)$ measuring gain over mixture $y$ (higher is better; Hungarian permutation matching offline). Vakyansh Hindi Wav2Vec2~\cite{chadha2021vakyansh} builds on wav2vec~2.0~\cite{baevski2020wav2vec} using CTC over Devanagari, while Stage~4 uses Pydantic schema validation and BM25 search~\cite{robertson2009bm25}.

\section{Method}
The pipeline comprises five sequential stages plus an extractive lexical query engine.

\noindent\textbf{Conditioning and separation:} Stage~0 standardizes audio to mono 16~kHz PCM, removes DC offset, normalizes amplitude to $-23$~dBFS RMS (peak ceiling $-1.0$~dBFS), and applies energy VAD (300~ms hangover). In Track~A, Stage~1 runs Conv-TasNet (\texttt{CONVTASNET\_BASE\_LIBRI2MIX}) using polyphase resampling (\texttt{scipy.signal.resample\_poly}, 16$\leftrightarrow$8~kHz), with Hybrid Demucs (\texttt{htdemucs}) as comparator. Track~B runs default pipeline settings (\texttt{htdemucs}/bypass) and does not evaluate Conv-TasNet.

\noindent\textbf{Diarization, ASR, and Structuring:} Stage~2 extracts sliding-window embeddings via a native 192-d acoustic-statistics hyperspherical projection extractor (Mel-filterbank mean/std/flux with QR orthogonal basis; \texttt{pyannote.audio} weights were absent in runtime) and clusters a cosine-affinity normalized random-walk Laplacian matrix via maximal eigengap $k^*$ and $k$-means (1.5~s/0.5~s windows in Track~A; 2.0~s/0.75~s in Track~B; 0.25~s NIST collar).

Stage~3 slices audio by diarization turns (50~ms padding) and transcribes via Vakyansh Hindi Wav2Vec2 (\texttt{Harveenchadha/}\allowbreak\texttt{vakyansh-wav2vec2-hindi-him-4200}) with greedy CTC argmax decoding (no language model rescoring), Unicode NFKC normalization, and nukta harmonization. An alternative comparator (\texttt{openai/whisper-tiny}) was evaluated on \texttt{mix\_01}.

Stage~4 structures transcripts into Pydantic~v2 schemas (\texttt{PostProcessedReport}, \texttt{SpeakerDialogue}, \texttt{KeyFinding}). Stage~4 quantitative metrics reflect offline deterministic template substitution via \texttt{MockLLMProvider}; separation, diarization, ASR, and profiling metrics come from their respective neural/signal modules. Under \texttt{MockLLMProvider}, 0.0\% hallucination reflects rule-based template fidelity rather than generative safety. Hallucination rate is $(|\mathcal{E}_{\text{unsupp}}| / |\mathcal{E}_{\text{total}}|) \times 100\%$ (lower is better; defaults to 0.0\% if zero candidate entities are extracted), matching summary entities/numbers against despaced CTC tokens and an Indic-English phonetic map. The \texttt{ConversationRAGEngine} retrieves turns via BM25 lexical scoring and synonym expansion without vector embeddings.

\section{Experimental setup}
\noindent\textbf{Datasets and Evaluation Tracks:} Track~A uses clean speech stems from 7 Hindi/Hinglish speakers (\texttt{speaker\_A}--\texttt{G}): repository documentation identifies Hugging Face \texttt{nameissakthi/hindi-english-bilingual} (\texttt{A}, \texttt{B}) and \texttt{Speech-data/Hindi-Speech-Dataset} (\texttt{D}), while \texttt{C}, \texttt{E}, \texttt{F}, and \texttt{G} lack upstream source metadata. Six 2-speaker mixtures test separation ($\Omega \in \{25\%, 50\%\}$, $\text{SIR} \in \{-6, 0, +6\}$~dB, 8.31--9.54~s); nine mixtures test diarization (adding 75\% overlap \texttt{mix\_07}, 3-speaker \texttt{mix\_08}, and 10~dB noise \texttt{mix\_09}); ASR is benchmarked on 9.54~s bilingual \texttt{controlled\_mix\_01} (pair~AD). Track~B evaluates three simulated domain-stress clips (27.3~s total) synthesized from Track~A stems in \texttt{scripts/expand\_data\_pipeline.py}: \texttt{air\_formal\_news\_01} (6.0~s, RMS-normalized \texttt{speaker\_D}, studio SNR 47.1~dB, 1 speaker), \texttt{air\_spontaneous\_panel\_01} (13.0~s, 3-speaker overlay of \texttt{A}, \texttt{F}, \texttt{G} with ambient noise, SNR 35.6~dB), and \texttt{vaani\_rural\_call\_01} (8.3~s, \texttt{speaker\_E} under 300--3400~Hz Butterworth bandpass + 100~Hz hum/noise, SNR 21.1~dB).

\noindent\textbf{Metrics and Environment:} Separation is scored by absolute SI-SDR and $\Delta$SI-SDR in dB (higher is better), diarization by NIST DER with 0.25~s collar (lower is better), ASR by normalized WER and CER (lower is better), and latency by $\text{RTF} = T_{\text{proc}} / T_{\text{audio}}$. Profiling ran on an 8-core x86\_64 WSL2 CPU host (\texttt{benchmark\_summary.json}), whereas earlier Stage~1 M2 and Track~B demo runs recorded Apple Silicon MPS device execution.

\section{Results}
\noindent\textbf{Controlled evaluations:} Table~\ref{tab:results} summarizes results across distinct controlled suites.

\begin{table}[ht]
\centering
\small
\begin{tabularx}{\linewidth}{@{}l l X@{}}
\toprule
Evaluation & Scope & Recorded result \\
\midrule
Separation & 6 mixtures & Conv-TasNet: mean $\Delta$SI-SDR +15.45~dB (peak +21.76~dB, min +4.81~dB, output 15.45~dB, baseline $-0.00$~dB; mean RTF 0.090). Demucs: output $-8.60$~dB ($\Delta$SI-SDR $-8.60$~dB; RTF 0.561). \\
Diarization & 9 mixtures & Mean DER: 40.52\% (raw spectral); 55.51\% (raw GMM); 23.97\% (with separation). Reduction: 16.55 pp. \\
ASR & 1 clip (9.54 s) & WER on \texttt{mix\_01}: 86.28\% (raw); 70.78\% (diarized); 51.57\% (separated+diarized; CER 44.66\%, RTF 0.194). Whisper: 100.0\% (raw); 103.34\% (separated). \\
\bottomrule
\end{tabularx}
\caption{Results across separate controlled evaluations (not directly comparable across rows).}
\label{tab:results}
\end{table}

Across the six two-speaker mixtures (\texttt{stage1\_separation\_results.json}), raw mixture SI-SDR averaged $-0.003$~dB ($-0.00$~dB rounded). Conv-TasNet reached 15.448~dB mean output SI-SDR, yielding $\Delta\text{SI-SDR} = +15.451$~dB ($+15.45$~dB; peak $+21.76$~dB on \texttt{mix\_01}, min $+4.81$~dB on \texttt{mix\_04}); committed Stage~1 JSON records mean RTF 0.090 (earlier M2 MPS report \texttt{e712501} logged 0.075). Hybrid Demucs averaged $-8.598$~dB output ($\Delta\text{SI-SDR} = -8.595$~dB, $-8.60$~dB rounded) due to music-stem grouping. Across nine mixtures, separation reduced mean DER from 40.52\% to 23.97\% (16.55 pp gain; GMM 55.51\%), with higher error on 3-speaker \texttt{mix\_08} (49.57\%) and 10~dB noisy \texttt{mix\_09} (46.60\%). On \texttt{controlled\_mix\_01}, Vakyansh WER fell from 86.28\% (raw) to 70.78\% (diarized) and 51.57\% (separated+diarized; CER 44.66\%, CPU RTF 0.194), whereas \texttt{whisper-tiny} suffered repetition collapse (100.0\% raw; 103.34\% separated).

In an exploratory analysis (\texttt{run\_ablation\_study.py}), seven predefined synthetic scenarios ($n=7$, \texttt{EXP-2.1}--\texttt{EXP-2.7}, $\Omega \in \{25\%, 40\%, 50\%, 75\%\}$) are configured separately from end-to-end execution (which uses \texttt{htdemucs} and \texttt{mock\_mode=True}). Across these synthetic scenarios, SI-SDR correlates inversely with DER ($r = -0.977$) and WER ($r = -0.969$), with configured boundary \texttt{EXP-2.5} (4.2~dB SI-SDR, 28.9\% DER, 42.7\% WER) marking an exploratory stress point rather than a universal empirical threshold.

\noindent\textbf{Simulated domain stress and profiling:} On Track~B: (1)~\texttt{air\_formal\_news\_01}: DER was 0.0\%; WER against the 41-word manifest was 92.7\% (8 subs, 30 deletions, CER 75.4\%) because the manifest combines altered Sentence~1 wording with an unuttered 24-word Sentence~2 (27 divergent/unuttered words vs.\ 30 edit-distance deletions from 11 glued tokens). Against the clip's actual 14 spoken words (\texttt{air\_formal\_news\_01\_diagnostic.json}), diagnostic WER is 57.1\% (8/14: 5 subs, 3 dels) with 0.0\% CER (69/69 chars match when AI is written in Devanagari), isolating CTC single-delimiter word merging. (2)~\texttt{air\_spontaneous\_panel\_01}: DER was 58.7\% and WER 57.8\% due to spectral cluster collapse ($k^* = 1$) under dense overlap. (3)~\texttt{vaani\_rural\_call\_01}: DER was 0.7\% and WER 63.0\% (CER 42.4\%) from 300--3400~Hz bandpass loss. On WSL2 CPU (\texttt{benchmark\_summary.json}), isolated 30~s end-to-end execution took 20.90~s (RTF 0.697; Stage~1 separation alone 19.40~s, RTF 0.647), continuous 30~s scaling achieved RTF 0.490, and cold-start Track~B news ran at RTF 19.98.

\section{Discussion and limitations}
Limitations include: (1)~\textbf{Fixed two-stream separation}: Conv-TasNet outputs two streams, leaving $N > 2$ talker overlap unresolved. (2)~\textbf{Domain mismatch}: \texttt{htdemucs} targets musical stems rather than multi-speaker separation. (3)~\textbf{Monolingual greedy CTC}: Vakyansh lacks bilingual LM rescoring, dropping Latin English words and merging adjacent Hindi words when single-space CTC delimiters fail. (4)~\textbf{Simulated Track~B corpus}: Field stress clips are synthetically constructed from Track~A stems (27.3~s total), serving as controlled stress diagnostics rather than large-scale wild audio. (5)~\textbf{Mock Stage~4 validation}: Schema contracts guarantee format compliance but rely on rule-based template grounding under \texttt{MockLLMProvider}.

\section{Conclusion}
We evaluated a modular five-stage pipeline for single-channel overlapping Hinglish speech across controlled mixtures and simulated broadcast/telephony stress clips. Pretrained Conv-TasNet improved separation fidelity ($\Delta$SI-SDR $+15.45$~dB over a $-0.00$~dB baseline), lowered mean DER from 40.52\% to 23.97\%, and reduced \texttt{controlled\_mix\_01} WER from 86.28\% to 51.57\%, whereas Stage~4 zero-hallucination scores reflect deterministic template validation. Stress diagnostics revealed spectral cluster collapse under dense overlap ($k^*=1$, 58.7\% DER), CTC delimiter word merging, and 300--3400~Hz bandpass degradation, motivating $N>2$ separation and bilingual LM decoding.

\section*{Individual member contributions}
{\small
\begin{tabularx}{\linewidth}{@{}l X@{}}
\toprule
\textbf{Name} & \textbf{Contribution} \\
\midrule
Kushal Gupta & Spectral clustering, Indic ASR normalizer, Pydantic schemas, CLI orchestration, error-propagation analysis, and RAG retrieval. \\
Vansh Sharma & Multi-speaker formulation, Conv-TasNet and Demucs integration, SI-SDR evaluation, speaker embeddings, and downstream ASR failure analysis. \\
Dorothy & Stage 0 audio preprocessing, VAD conditioning, spectrogram analysis, and curation of evaluation clips and manifests. \\
Pranav & Pydantic schema contracts, benchmark profiling, latency and RTF benchmarking, and pipeline test coverage. \\
\bottomrule
\end{tabularx}}

\smallskip
\noindent All members have read this report and agree with the statement above.

\section*{Generative AI disclosure}
AI assistants were used for targeted engineering and editing support: Anthropic Claude Opus 4.6 assisted in PyTorch polyphase resampling and Stage~1 tensor-dimension debugging; Google Gemini 3.8 Flash assisted in Stage~4 Pydantic~v2 schema and prompt-template drafting; OpenAI GPT-4o (accessed via ChatGPT Plus) aided LaTeX table formatting and proofreading; and the Google Antigravity agent (powered by Google Gemini 3.8 Flash and Gemini 4 models) assisted in code consistency audits, test validation, and report/README/demo synchronization. All experimental protocols, code, benchmarks, and analyses were executed and verified by the student authors, who take full responsibility for this submission.

\section*{Code availability}
\noindent Source code: \url{\githuburl} (public, includes reproduction scripts and 127 passing unit tests). Interactive Web Demo: \url{\demourl}

\begin{thebibliography}{11}
\small
\bibitem{luo2019convtasnet}
Y.~Luo and N.~Mesgarani.
\newblock Conv-TasNet: Surpassing ideal time--frequency magnitude masking for speech separation.
\newblock \emph{IEEE/ACM Trans. Audio, Speech, Lang. Process.}, 27(8):1256--1266, 2019.

\bibitem{defossez2021demucs}
A.~D{\'e}fossez.
\newblock Hybrid spectrogram and waveform source separation.
\newblock In \emph{Proc. ISMIR Workshop on Music Separation with DNNs}, 2021.

\bibitem{bredin2020pyannote}
H.~Bredin, R.~Yin, J.~M. Coria, G.~Broux, et~al.
\newblock pyannote.audio: neural building blocks for speaker diarization.
\newblock In \emph{Proc. ICASSP}, pages 7124--7128, 2020.

\bibitem{baevski2020wav2vec}
A.~Baevski, H.~Zhou, A.~Mohamed, and M.~Auli.
\newblock wav2vec 2.0: A framework for self-supervised learning of speech representations.
\newblock In \emph{Advances in Neural Information Processing Systems}, volume~33, pages 12449--12460, 2020.

\bibitem{chadha2021vakyansh}
H.~Chadha.
\newblock Vakyansh Hindi Wav2Vec2 CTC model.
\newblock Hugging Face model card, \url{https://huggingface.co/Harveenchadha/vakyansh-wav2vec2-hindi-him-4200}, 2021.

\bibitem{cosentino2020librimix}
J.~Cosentino, M.~Pariente, D.~Wang, and A.~Deleforge.
\newblock LibriMix: An open-source dataset for general sound source separation.
\newblock \emph{arXiv preprint arXiv:2005.11262}, 2020.

\bibitem{yang2022torchaudio}
Y.-Y. Yang, M.~Hira, Z.~Ni, A.~Astafurov, et~al.
\newblock Torchaudio: Building blocks for audio and speech processing.
\newblock In \emph{Proc. ICASSP}, pages 6982--6986, 2022.

\bibitem{leroux2019sisdr}
J.~Le~Roux, S.~Wisdom, H.~Erdogan, and J.~R. Hershey.
\newblock SDR -- half-baked or well done?
\newblock In \emph{Proc. ICASSP}, pages 626--630, 2019.

\bibitem{vincent2006performance}
E.~Vincent, R.~Gribonval, and C.~F{\'e}votte.
\newblock Performance measurement in blind audio source separation.
\newblock \emph{IEEE Trans. Audio, Speech, Lang. Process.}, 14(4):1462--1469, 2006.

\bibitem{ng2001spectral}
A.~Y. Ng, M.~I. Jordan, and Y.~Weiss.
\newblock On spectral clustering: Analysis and an algorithm.
\newblock In \emph{Advances in Neural Information Processing Systems}, volume~14, pages 849--856, 2001.

\bibitem{robertson2009bm25}
S.~Robertson and H.~Zaragoza.
\newblock The probabilistic relevance framework: BM25 and beyond.
\newblock \emph{Foundations and Trends in Information Retrieval}, 3(4):333--389, 2009.
\end{thebibliography}

\end{document}
